% Source: source/普通高考/2013/2013湖北理.pdf, page 1, question 12.
% Topology: start -> a=10,i=1 -> a=4?; 是 -> 输出 i -> 结束;
% 否 -> a 是奇数?; 是 -> a=3a+1, 否 -> a=a/2; both -> i=i+1 -> loop to a=4?.
\begin{tikzpicture}[
  x=1cm, y=0.92cm, >=Stealth,
  line width=0.45pt, line cap=round, line join=round,
  every node/.style={anchor=center, align=center,
    execute at begin node={\setlength{\parindent}{0pt}\noindent}},
  flowbox/.style={draw, minimum height=0.70cm, inner xsep=4pt, inner ysep=2pt,
    align=center, execute at begin node={\setlength{\parindent}{0pt}\noindent}},
  terminal/.style={flowbox, rounded corners=0.16cm, minimum width=1.30cm,
    minimum height=0.72cm},
  process/.style={flowbox, minimum width=2.75cm},
  decision/.style={flowbox, diamond, aspect=3.0, minimum width=3.75cm,
    minimum height=1.00cm, inner sep=1pt},
  io/.style={flowbox, trapezium, trapezium left angle=72, trapezium right angle=108,
    minimum width=1.95cm, minimum height=0.76cm},
  branchlabel/.style={anchor=center, align=center, inner sep=0pt,
    execute at begin node={\setlength{\parindent}{0pt}\noindent}}
]
  \node[terminal] (start) at (0,10.0) {开始};
  \node[process, minimum width=3.35cm] (init) at (0,8.65) {$a=10,\ i=1$};
  \node[decision] (testA4) at (0,7.05) {$a=4?$};
  \node[decision, minimum width=4.35cm] (testOdd) at (0,5.15) {$a\text{ 是奇数?}$};
  \node[process, minimum width=3.05cm] (odd) at (-4.45,3.35) {$a=3a+1$};
  \node[process, minimum width=2.55cm] (half) at (3.95,3.35) {$a=\dfrac{a}{2}$};
  \node[process, minimum width=2.55cm] (inc) at (0,1.15) {$i=i+1$};
  \node[io] (output) at (6.15,6.05) {输出 $i$};
  \node[terminal] (finish) at (6.15,4.55) {结束};

  \draw[->] (start.south) -- (init.north);
  \coordinate (loopJoin) at (0,7.70);
  \draw (init.south) -- (loopJoin);
  \draw[->] (loopJoin) -- (testA4.north);
  \draw[->] (testA4.south) -- node[branchlabel, right=2pt] {否} (testOdd.north);
  \draw[->] (testA4.east) -- node[branchlabel, above=2pt] {是} (6.15,7.05) -- (output.north);
  \draw[->] (output.south) -- (finish.north);

  \draw (testOdd.west) -- node[branchlabel, above=2pt] {是} (-4.45,5.15);
  \draw[->] (-4.45,5.15) -- (odd.north);
  \draw (odd.south) -- (-4.45,2.05) -- (0,2.05);
  \draw[->] (testOdd.east) -- node[branchlabel, above=2pt] {否} (3.95,5.15) -- (half.north);
  \draw (half.south) -- (3.95,2.05);
  \draw[->] (3.95,2.05) -- (0,2.05) -- (inc.north);

  \coordinate (loopDrop) at (0,0.35);
  \coordinate (loopLeft) at (-7.05,0.35);
  \coordinate (loopTop) at (-7.05,7.70);
  \draw (inc.south) -- (loopDrop) -- (loopLeft) -- (loopTop);
  \draw[->] (loopTop) -- (loopJoin);
\end{tikzpicture}
